%%%%%%%%%%%%%%%%%%%%%%%%%%%%%%%%%%%%%%%%%%%%%%%%%%%%%%%%%%%%%%%%%%%%%%%%
% Section conclusion
% Inclus depuis main.tex via %%%%%%%%%%%%%%%%%%%%%%%%%%%%%%%%%%%%%%%%%%%%%%%%%%%%%%%%%%%%%%%%%%%%%%%%
% Section conclusion
% Inclus depuis main.tex via %%%%%%%%%%%%%%%%%%%%%%%%%%%%%%%%%%%%%%%%%%%%%%%%%%%%%%%%%%%%%%%%%%%%%%%%
% Section conclusion
% Inclus depuis main.tex via %%%%%%%%%%%%%%%%%%%%%%%%%%%%%%%%%%%%%%%%%%%%%%%%%%%%%%%%%%%%%%%%%%%%%%%%
% Section conclusion
% Inclus depuis main.tex via \input{chapters/conclusion}
%%%%%%%%%%%%%%%%%%%%%%%%%%%%%%%%%%%%%%%%%%%%%%%%%%%%%%%%%%%%%%%%%%%%%%%%

%%%%%%%%%%%%%%%%%%%%%%%%%%%%%%%%%%%%%%%%%%%%%%%%%%%%%%%%%%%%%%%%%%%%%%%%

\subsection{Réponses aux questions de recherche}

Les travaux présentés apportent des réponses nuancées aux quatre questions posées en
introduction. Concernant la génération (Q1), les grands modèles de langue peuvent
effectivement produire des documents cliniques linguistiquement réalistes et médicalement
cohérents, à condition d'injecter la connaissance du domaine via une ontologie
(extraction de concepts UMLS) et de raffiner itérativement les sorties à partir d'un
signal de récompense issu d'un scoreur contrastif. L'approche fac-similé, en préservant
une correspondance avec des narratifs authentiques, capture la structure documentaire et
le raisonnement clinique que la génération \textit{ex nihilo} peine à reproduire.

Sur l'utilité aval (Q2), les documents synthétiques se révèlent efficaces comme complément
plutôt que substitut des données réelles. Les gains sont maximaux dans les régimes à faibles
ressources : classes minoritaires, pathologies rares, contextes institutionnels où l'annotation
manuelle est prohibitive. La complexité de la tâche module l'ampleur du bénéfice : les
tâches de classification fine (CIM-9 top-100, CIM-10 top-100) bénéficient davantage de
l'alignement itératif que les tâches plus simples. Sur le français, le scoreur composite
(BioLORD + LLM-juge) produit des gains relatifs allant jusqu'à \SI{37}{\percent} sur les
labels rares, confirmant l'importance du signal de récompense.

Concernant l'annotation faible (Q3), les grands modèles de langue constituent des
annotateurs compétents, particulièrement précieux pour les langues faiblement dotées
comme le basque. Leurs erreurs sont complémentaires de celles des approches
dictionnaires : recall large mais sur-extraction d'un côté, précision étroite de l'autre. La
stratégie de supervision hybride, qui mélange ces deux sources d'annotations à un ratio
optimisé par langue, surpasse chacune prise isolément. La distillation vers des modèles
encodeurs locaux préserve la souveraineté des données et s'avère souvent supérieure à
l'inférence directe du grand modèle, grâce au filtrage implicite des erreurs LLM.

Enfin, sur la vie privée (Q4), la thèse apporte un constat lucide : la génération synthétique
\textit{réduit} mais n'\textit{élimine} pas les risques de réidentification. Les attaques de
discrimination atteignent \SIrange{89}{96}{\percent} de succès sur les documents de base,
et les attaques par liaison \SIrange{45}{52}{\percent}, bien au-dessus du hasard. La cause
identifiée est la combinaison distinctive de mots-clés UMLS qui encode une empreinte
patient. De manière contre-intuitive, l'alignement DPO amplifie ce risque par rapport au
SFT seul. La perturbation ciblée des mots-clés (remplacement par synonymes d'ontologie)
réduit le risque de liaison de \SI{37}{\percent} absolus pour seulement \SI{3}{\percent} de
perte d'utilité aval, offrant un compromis déployable en pratique.

\subsection{Limites et arbitrage utilité/confidentialité}

Aucune solution parfaite n'émerge : résiduellement, des risques demeurent. Plusieurs limites
méritent d'être soulignées. D'une part, les annotations dites \emph{silver} issues du pipeline
d'annotation faible imposent un plafond empirique de performance. D'autre part, la validation
de l'approche a porté sur des notes cliniques généralistes et cinq langues européennes :
les spécialités (radiologie, oncologie, anatomopathologie) et les langues faiblement dotées
hors Europe restent à explorer. Enfin, l'approche fournit des garanties empiriques par
évaluation d'attaques, non des garanties formelles à la manière de la confidentialité
différentielle. L'articulation entre garanties formelles et préservation de l'utilité constitue
une question ouverte.

\subsection{Perspectives}

Plusieurs directions prolongent naturellement ces travaux. Sur le plan technique, l'extension
du fac-similé à des dossiers patients multimodaux longitudinaux, combinant texte, données
structurées (constantes, analyses biologiques) et séries temporelles, permettrait de répondre
à des cas d'usage au-delà du seul TAL : aide à la décision, formation médicale, simulation de
parcours de soins. Le découplage complet vis-à-vis des documents d'amorçage, par exemple
en tirant des séquences de pseudo-mots-clés d'un petit modèle entraîné sur les co-occurrences
de l'ontologie médicale, est déjà amorcé dans le volet français et mérite d'être consolidé.

Sur le plan méthodologique, la combinaison de la confidentialité différentielle et du
raffinement par préférence, ainsi que la quantification théorique de la relation entre taille
du corpus synthétique et risque de réidentification, offriraient aux institutions un cadre de
calibration explicite. Enfin, les approches développées ne sont pas spécifiques au domaine
médical : partout où des textes sensibles doivent être exploités sans exposer d'information
personnelle (droit, finance, assurance), le même arbitrage utilité/confidentialité se pose et
les mêmes stratégies fac-similé et de supervision faible peuvent s'appliquer.

La vision générale est celle d'une \textit{innovation responsable} : faire progresser les
capacités du TAL clinique en respectant la vie privée des patients, en autonomisant les
institutions de santé et en démocratisant l'accès à des outils qui améliorent la qualité des
soins et l'efficacité de la recherche. Atteindre cet objectif suppose un travail technique
continu mais, tout autant, la construction de normes, de standards et de cadres de
gouvernance permettant le déploiement sûr et équitable des données synthétiques et des
méthodes fondées sur les grands modèles de langue. Cette thèse contribue à cet horizon
en proposant des méthodes, en quantifiant leurs limites et en assumant l'absence de
solution à risque nul.

%%%%%%%%%%%%%%%%%%%%%%%%%%%%%%%%%%%%%%%%%%%%%%%%%%%%%%%%%%%%%%%%%%%%%%%%

%%%%%%%%%%%%%%%%%%%%%%%%%%%%%%%%%%%%%%%%%%%%%%%%%%%%%%%%%%%%%%%%%%%%%%%%

\subsection{Réponses aux questions de recherche}

Les travaux présentés apportent des réponses nuancées aux quatre questions posées en
introduction. Concernant la génération (Q1), les grands modèles de langue peuvent
effectivement produire des documents cliniques linguistiquement réalistes et médicalement
cohérents, à condition d'injecter la connaissance du domaine via une ontologie
(extraction de concepts UMLS) et de raffiner itérativement les sorties à partir d'un
signal de récompense issu d'un scoreur contrastif. L'approche fac-similé, en préservant
une correspondance avec des narratifs authentiques, capture la structure documentaire et
le raisonnement clinique que la génération \textit{ex nihilo} peine à reproduire.

Sur l'utilité aval (Q2), les documents synthétiques se révèlent efficaces comme complément
plutôt que substitut des données réelles. Les gains sont maximaux dans les régimes à faibles
ressources : classes minoritaires, pathologies rares, contextes institutionnels où l'annotation
manuelle est prohibitive. La complexité de la tâche module l'ampleur du bénéfice : les
tâches de classification fine (CIM-9 top-100, CIM-10 top-100) bénéficient davantage de
l'alignement itératif que les tâches plus simples. Sur le français, le scoreur composite
(BioLORD + LLM-juge) produit des gains relatifs allant jusqu'à \SI{37}{\percent} sur les
labels rares, confirmant l'importance du signal de récompense.

Concernant l'annotation faible (Q3), les grands modèles de langue constituent des
annotateurs compétents, particulièrement précieux pour les langues faiblement dotées
comme le basque. Leurs erreurs sont complémentaires de celles des approches
dictionnaires : recall large mais sur-extraction d'un côté, précision étroite de l'autre. La
stratégie de supervision hybride, qui mélange ces deux sources d'annotations à un ratio
optimisé par langue, surpasse chacune prise isolément. La distillation vers des modèles
encodeurs locaux préserve la souveraineté des données et s'avère souvent supérieure à
l'inférence directe du grand modèle, grâce au filtrage implicite des erreurs LLM.

Enfin, sur la vie privée (Q4), la thèse apporte un constat lucide : la génération synthétique
\textit{réduit} mais n'\textit{élimine} pas les risques de réidentification. Les attaques de
discrimination atteignent \SIrange{89}{96}{\percent} de succès sur les documents de base,
et les attaques par liaison \SIrange{45}{52}{\percent}, bien au-dessus du hasard. La cause
identifiée est la combinaison distinctive de mots-clés UMLS qui encode une empreinte
patient. De manière contre-intuitive, l'alignement DPO amplifie ce risque par rapport au
SFT seul. La perturbation ciblée des mots-clés (remplacement par synonymes d'ontologie)
réduit le risque de liaison de \SI{37}{\percent} absolus pour seulement \SI{3}{\percent} de
perte d'utilité aval, offrant un compromis déployable en pratique.

\subsection{Limites et arbitrage utilité/confidentialité}

Aucune solution parfaite n'émerge : résiduellement, des risques demeurent. Plusieurs limites
méritent d'être soulignées. D'une part, les annotations dites \emph{silver} issues du pipeline
d'annotation faible imposent un plafond empirique de performance. D'autre part, la validation
de l'approche a porté sur des notes cliniques généralistes et cinq langues européennes :
les spécialités (radiologie, oncologie, anatomopathologie) et les langues faiblement dotées
hors Europe restent à explorer. Enfin, l'approche fournit des garanties empiriques par
évaluation d'attaques, non des garanties formelles à la manière de la confidentialité
différentielle. L'articulation entre garanties formelles et préservation de l'utilité constitue
une question ouverte.

\subsection{Perspectives}

Plusieurs directions prolongent naturellement ces travaux. Sur le plan technique, l'extension
du fac-similé à des dossiers patients multimodaux longitudinaux, combinant texte, données
structurées (constantes, analyses biologiques) et séries temporelles, permettrait de répondre
à des cas d'usage au-delà du seul TAL : aide à la décision, formation médicale, simulation de
parcours de soins. Le découplage complet vis-à-vis des documents d'amorçage, par exemple
en tirant des séquences de pseudo-mots-clés d'un petit modèle entraîné sur les co-occurrences
de l'ontologie médicale, est déjà amorcé dans le volet français et mérite d'être consolidé.

Sur le plan méthodologique, la combinaison de la confidentialité différentielle et du
raffinement par préférence, ainsi que la quantification théorique de la relation entre taille
du corpus synthétique et risque de réidentification, offriraient aux institutions un cadre de
calibration explicite. Enfin, les approches développées ne sont pas spécifiques au domaine
médical : partout où des textes sensibles doivent être exploités sans exposer d'information
personnelle (droit, finance, assurance), le même arbitrage utilité/confidentialité se pose et
les mêmes stratégies fac-similé et de supervision faible peuvent s'appliquer.

La vision générale est celle d'une \textit{innovation responsable} : faire progresser les
capacités du TAL clinique en respectant la vie privée des patients, en autonomisant les
institutions de santé et en démocratisant l'accès à des outils qui améliorent la qualité des
soins et l'efficacité de la recherche. Atteindre cet objectif suppose un travail technique
continu mais, tout autant, la construction de normes, de standards et de cadres de
gouvernance permettant le déploiement sûr et équitable des données synthétiques et des
méthodes fondées sur les grands modèles de langue. Cette thèse contribue à cet horizon
en proposant des méthodes, en quantifiant leurs limites et en assumant l'absence de
solution à risque nul.

%%%%%%%%%%%%%%%%%%%%%%%%%%%%%%%%%%%%%%%%%%%%%%%%%%%%%%%%%%%%%%%%%%%%%%%%

%%%%%%%%%%%%%%%%%%%%%%%%%%%%%%%%%%%%%%%%%%%%%%%%%%%%%%%%%%%%%%%%%%%%%%%%

\subsection{Réponses aux questions de recherche}

Les travaux présentés apportent des réponses nuancées aux quatre questions posées en
introduction. Concernant la génération (Q1), les grands modèles de langue peuvent
effectivement produire des documents cliniques linguistiquement réalistes et médicalement
cohérents, à condition d'injecter la connaissance du domaine via une ontologie
(extraction de concepts UMLS) et de raffiner itérativement les sorties à partir d'un
signal de récompense issu d'un scoreur contrastif. L'approche fac-similé, en préservant
une correspondance avec des narratifs authentiques, capture la structure documentaire et
le raisonnement clinique que la génération \textit{ex nihilo} peine à reproduire.

Sur l'utilité aval (Q2), les documents synthétiques se révèlent efficaces comme complément
plutôt que substitut des données réelles. Les gains sont maximaux dans les régimes à faibles
ressources : classes minoritaires, pathologies rares, contextes institutionnels où l'annotation
manuelle est prohibitive. La complexité de la tâche module l'ampleur du bénéfice : les
tâches de classification fine (CIM-9 top-100, CIM-10 top-100) bénéficient davantage de
l'alignement itératif que les tâches plus simples. Sur le français, le scoreur composite
(BioLORD + LLM-juge) produit des gains relatifs allant jusqu'à \SI{37}{\percent} sur les
labels rares, confirmant l'importance du signal de récompense.

Concernant l'annotation faible (Q3), les grands modèles de langue constituent des
annotateurs compétents, particulièrement précieux pour les langues faiblement dotées
comme le basque. Leurs erreurs sont complémentaires de celles des approches
dictionnaires : recall large mais sur-extraction d'un côté, précision étroite de l'autre. La
stratégie de supervision hybride, qui mélange ces deux sources d'annotations à un ratio
optimisé par langue, surpasse chacune prise isolément. La distillation vers des modèles
encodeurs locaux préserve la souveraineté des données et s'avère souvent supérieure à
l'inférence directe du grand modèle, grâce au filtrage implicite des erreurs LLM.

Enfin, sur la vie privée (Q4), la thèse apporte un constat lucide : la génération synthétique
\textit{réduit} mais n'\textit{élimine} pas les risques de réidentification. Les attaques de
discrimination atteignent \SIrange{89}{96}{\percent} de succès sur les documents de base,
et les attaques par liaison \SIrange{45}{52}{\percent}, bien au-dessus du hasard. La cause
identifiée est la combinaison distinctive de mots-clés UMLS qui encode une empreinte
patient. De manière contre-intuitive, l'alignement DPO amplifie ce risque par rapport au
SFT seul. La perturbation ciblée des mots-clés (remplacement par synonymes d'ontologie)
réduit le risque de liaison de \SI{37}{\percent} absolus pour seulement \SI{3}{\percent} de
perte d'utilité aval, offrant un compromis déployable en pratique.

\subsection{Limites et arbitrage utilité/confidentialité}

Aucune solution parfaite n'émerge : résiduellement, des risques demeurent. Plusieurs limites
méritent d'être soulignées. D'une part, les annotations dites \emph{silver} issues du pipeline
d'annotation faible imposent un plafond empirique de performance. D'autre part, la validation
de l'approche a porté sur des notes cliniques généralistes et cinq langues européennes :
les spécialités (radiologie, oncologie, anatomopathologie) et les langues faiblement dotées
hors Europe restent à explorer. Enfin, l'approche fournit des garanties empiriques par
évaluation d'attaques, non des garanties formelles à la manière de la confidentialité
différentielle. L'articulation entre garanties formelles et préservation de l'utilité constitue
une question ouverte.

\subsection{Perspectives}

Plusieurs directions prolongent naturellement ces travaux. Sur le plan technique, l'extension
du fac-similé à des dossiers patients multimodaux longitudinaux, combinant texte, données
structurées (constantes, analyses biologiques) et séries temporelles, permettrait de répondre
à des cas d'usage au-delà du seul TAL : aide à la décision, formation médicale, simulation de
parcours de soins. Le découplage complet vis-à-vis des documents d'amorçage, par exemple
en tirant des séquences de pseudo-mots-clés d'un petit modèle entraîné sur les co-occurrences
de l'ontologie médicale, est déjà amorcé dans le volet français et mérite d'être consolidé.

Sur le plan méthodologique, la combinaison de la confidentialité différentielle et du
raffinement par préférence, ainsi que la quantification théorique de la relation entre taille
du corpus synthétique et risque de réidentification, offriraient aux institutions un cadre de
calibration explicite. Enfin, les approches développées ne sont pas spécifiques au domaine
médical : partout où des textes sensibles doivent être exploités sans exposer d'information
personnelle (droit, finance, assurance), le même arbitrage utilité/confidentialité se pose et
les mêmes stratégies fac-similé et de supervision faible peuvent s'appliquer.

La vision générale est celle d'une \textit{innovation responsable} : faire progresser les
capacités du TAL clinique en respectant la vie privée des patients, en autonomisant les
institutions de santé et en démocratisant l'accès à des outils qui améliorent la qualité des
soins et l'efficacité de la recherche. Atteindre cet objectif suppose un travail technique
continu mais, tout autant, la construction de normes, de standards et de cadres de
gouvernance permettant le déploiement sûr et équitable des données synthétiques et des
méthodes fondées sur les grands modèles de langue. Cette thèse contribue à cet horizon
en proposant des méthodes, en quantifiant leurs limites et en assumant l'absence de
solution à risque nul.

%%%%%%%%%%%%%%%%%%%%%%%%%%%%%%%%%%%%%%%%%%%%%%%%%%%%%%%%%%%%%%%%%%%%%%%%

%%%%%%%%%%%%%%%%%%%%%%%%%%%%%%%%%%%%%%%%%%%%%%%%%%%%%%%%%%%%%%%%%%%%%%%%

\subsection{Réponses aux questions de recherche}

Les travaux présentés apportent des réponses nuancées aux quatre questions posées en
introduction. Concernant la génération (Q1), les grands modèles de langue peuvent
effectivement produire des documents cliniques linguistiquement réalistes et médicalement
cohérents, à condition d'injecter la connaissance du domaine via une ontologie
(extraction de concepts UMLS) et de raffiner itérativement les sorties à partir d'un
signal de récompense issu d'un scoreur contrastif. L'approche fac-similé, en préservant
une correspondance avec des narratifs authentiques, capture la structure documentaire et
le raisonnement clinique que la génération \textit{ex nihilo} peine à reproduire.

Sur l'utilité aval (Q2), les documents synthétiques se révèlent efficaces comme complément
plutôt que substitut des données réelles. Les gains sont maximaux dans les régimes à faibles
ressources : classes minoritaires, pathologies rares, contextes institutionnels où l'annotation
manuelle est prohibitive. La complexité de la tâche module l'ampleur du bénéfice : les
tâches de classification fine (CIM-9 top-100, CIM-10 top-100) bénéficient davantage de
l'alignement itératif que les tâches plus simples. Sur le français, le scoreur composite
(BioLORD + LLM-juge) produit des gains relatifs allant jusqu'à \SI{37}{\percent} sur les
labels rares, confirmant l'importance du signal de récompense.

Concernant l'annotation faible (Q3), les grands modèles de langue constituent des
annotateurs compétents, particulièrement précieux pour les langues faiblement dotées
comme le basque. Leurs erreurs sont complémentaires de celles des approches
dictionnaires : recall large mais sur-extraction d'un côté, précision étroite de l'autre. La
stratégie de supervision hybride, qui mélange ces deux sources d'annotations à un ratio
optimisé par langue, surpasse chacune prise isolément. La distillation vers des modèles
encodeurs locaux préserve la souveraineté des données et s'avère souvent supérieure à
l'inférence directe du grand modèle, grâce au filtrage implicite des erreurs LLM.

Enfin, sur la vie privée (Q4), la thèse apporte un constat lucide : la génération synthétique
\textit{réduit} mais n'\textit{élimine} pas les risques de réidentification. Les attaques de
discrimination atteignent \SIrange{89}{96}{\percent} de succès sur les documents de base,
et les attaques par liaison \SIrange{45}{52}{\percent}, bien au-dessus du hasard. La cause
identifiée est la combinaison distinctive de mots-clés UMLS qui encode une empreinte
patient. De manière contre-intuitive, l'alignement DPO amplifie ce risque par rapport au
SFT seul. La perturbation ciblée des mots-clés (remplacement par synonymes d'ontologie)
réduit le risque de liaison de \SI{37}{\percent} absolus pour seulement \SI{3}{\percent} de
perte d'utilité aval, offrant un compromis déployable en pratique.

\subsection{Limites et arbitrage utilité/confidentialité}

Aucune solution parfaite n'émerge : résiduellement, des risques demeurent. Plusieurs limites
méritent d'être soulignées. D'une part, les annotations dites \emph{silver} issues du pipeline
d'annotation faible imposent un plafond empirique de performance. D'autre part, la validation
de l'approche a porté sur des notes cliniques généralistes et cinq langues européennes :
les spécialités (radiologie, oncologie, anatomopathologie) et les langues faiblement dotées
hors Europe restent à explorer. Enfin, l'approche fournit des garanties empiriques par
évaluation d'attaques, non des garanties formelles à la manière de la confidentialité
différentielle. L'articulation entre garanties formelles et préservation de l'utilité constitue
une question ouverte.

\subsection{Perspectives}

Plusieurs directions prolongent naturellement ces travaux. Sur le plan technique, l'extension
du fac-similé à des dossiers patients multimodaux longitudinaux, combinant texte, données
structurées (constantes, analyses biologiques) et séries temporelles, permettrait de répondre
à des cas d'usage au-delà du seul TAL : aide à la décision, formation médicale, simulation de
parcours de soins. Le découplage complet vis-à-vis des documents d'amorçage, par exemple
en tirant des séquences de pseudo-mots-clés d'un petit modèle entraîné sur les co-occurrences
de l'ontologie médicale, est déjà amorcé dans le volet français et mérite d'être consolidé.

Sur le plan méthodologique, la combinaison de la confidentialité différentielle et du
raffinement par préférence, ainsi que la quantification théorique de la relation entre taille
du corpus synthétique et risque de réidentification, offriraient aux institutions un cadre de
calibration explicite. Enfin, les approches développées ne sont pas spécifiques au domaine
médical : partout où des textes sensibles doivent être exploités sans exposer d'information
personnelle (droit, finance, assurance), le même arbitrage utilité/confidentialité se pose et
les mêmes stratégies fac-similé et de supervision faible peuvent s'appliquer.

La vision générale est celle d'une \textit{innovation responsable} : faire progresser les
capacités du TAL clinique en respectant la vie privée des patients, en autonomisant les
institutions de santé et en démocratisant l'accès à des outils qui améliorent la qualité des
soins et l'efficacité de la recherche. Atteindre cet objectif suppose un travail technique
continu mais, tout autant, la construction de normes, de standards et de cadres de
gouvernance permettant le déploiement sûr et équitable des données synthétiques et des
méthodes fondées sur les grands modèles de langue. Cette thèse contribue à cet horizon
en proposant des méthodes, en quantifiant leurs limites et en assumant l'absence de
solution à risque nul.
